\subsection{Continuous Treatment Effect Bounds Without Ignorability}
\label{sec:bounds}
The CAPO and APO (dose-response) functions cannot be point identified from observational data without ignorability.
Under the CMSM with a given $\Lambda$, we can only identify a set of CAPO and APO functions jointly consistent with the observational data $\D$ and the continuous marginal sensitivity model.
All of the functions in this set are possible from the point of view of the observational data alone. 
So to cover the range of all possible functional values, we seek an interval function that maps covariate values, $\X=\x$, to the upper and lower bounds of this set for every treatment value, $\tf$.

For $\tf \in \Tcal$ and $\x \in \Xcal$, let $p(\yt \mid \tf, \x)$ denote the density of the distribution $P(\Yt \mid \Tf=\tf, \X=\x)$.
As a reminder, this distribution is identifiable from observational data, but without further assumptions the CAPO, $\mu(\x, \tf) = \E\left[ \Yt \mid \X = \x \right]$, is not.
We can express the CAPO in terms of its identifiable and unidentifiable components as
\begin{equation}
    \begin{split}
        \mu(\x, \tf)
        &= \frac{\int_\Ycal \yt \frac{p(\yt \mid \tf, \x)}{p(\tf \mid \yt, \x)}d\yt}{\int_\Ycal \frac{p(\yt \mid \tf, \x)}{p(\tf \mid \yt, \x)} d\yt}
        = \mut(\x, \tf) + \frac{\int_\Ycal w(\y, \x) (\y - \mut(\x, \tf)) p(\y \mid \tf, \x)d\y}{(\Lambda^2 - 1)^{-1} + \int_\Ycal w(\y, \x)p(\y \mid \tf, \x)d\y}, \\
        &\equiv \mu(w(\y, \x); \x, \tf, \Lambda)
    \end{split}
    \label{eq:mu_w}
\end{equation}
where, by a one-to-one change of variables, $\frac{1}{p(\tf \mid \yt, \x)} = \frac{1}{\Lambda p(\tf \mid \x)} + w(\y, \x) (\frac{\Lambda}{p(\tf \mid \x)} - \frac{1}{\Lambda p(\tf \mid \x)})$ with $w : \Ycal\times \Xcal \to [0, 1]$.
Both \cite{kallus2019interval} and later \cite{jesson2021quantifying} provide analogous expressions for the CAPO in the discrete treatment regime under the MSM, and we provide our derivation in \Cref{lem:capo}. 

The uncertainty set that includes all possible values of $w(\y, \x)$ that agree with the CMSM, \emph{i.e.}, the set of functions that violate ignorability by no more than $\Lambda$, can now be expressed as 
$
% \begin{equation*}
    \mathcal{W} = \{w: w(\y, \x) \in [ 0, 1 ] \quad \forall \y \in \Ycal,  \forall \x \in \Xcal \}.
$
% \end{equation*}

With this set of functions, we can now define the CAPO and APO bounds under the CMSM.
The CAPO lower, $\underline{\mu}(\x, \tf; \Lambda)$, and upper, $\overline{\mu}(\x, \tf; \Lambda)$, bounds under the CMSM with parameter $\Lambda$ are:
\noindent
\begin{minipage}{.5\linewidth}
    \begin{equation}
        \label{eq:CAPO_lower}
        \begin{split}
            \underline{\mu}(\x, \tf; \Lambda) 
            &\coloneqq \inf_{w \in \mathcal{W}} \mu(w(\y, \x); \x, \tf, \Lambda) \\
            &= \inf_{w \in \WniH} \mu(w(\y); \x, \tf, \Lambda) 
        \end{split}
    \end{equation}
\end{minipage}%
\begin{minipage}{.5\linewidth}
    \begin{equation}
        \label{eq:CAPO_upper}
        \begin{split}
            \overline{\mu}(\x, \tf; \Lambda) 
            &\coloneqq \sup_{w \in \mathcal{W}} \mu(w(\y, \x); \x, \tf, \Lambda) \\
            &= \sup_{w \in \WndH} \mu(w(\y); \x, \tf, \Lambda)
        \end{split}
    \end{equation}
\end{minipage}
\break
Where the sets $\WniH = \left\{ w: w(y) = H(\yH - \y) \right\}_{\yH \in \Ycal}$, and $\WndH = \left\{ w: w(y) = H(\y - \yH)  \right\}_{\yH \in \Ycal}$, and $H(\cdot)$ is the Heaviside step function. 
\Cref{lem:monotonic} in \cref{app:step} proves the equivalence in \cref{eq:CAPO_upper} for bounded $Y$. 
The equivalence in \cref{eq:CAPO_lower} can be proved analogously.

The APO lower, $\underline{\mu}(\tf; \Lambda)$, and upper, $\overline{\mu}(\tf; \Lambda)$, bounds under the CMSM with parameter $\Lambda$ are:
\noindent
\begin{minipage}{.5\linewidth}
    \begin{equation}
        \label{eq:APO_lower}
        \underline{\mu}(\tf; \Lambda) \coloneqq \E\left[ \underline{\mu}(\X, \tf; \Lambda) \right]
    \end{equation}
\end{minipage}%
\begin{minipage}{.5\linewidth}
    \begin{equation}
        \label{eq:APO_upper}
        \overline{\mu}(\tf; \Lambda) \coloneqq \E\left[ \overline{\mu}(\X, \tf; \Lambda) \right]
    \end{equation}
\end{minipage}

\textbf{Remark.} It is worth pausing here and breaking down \Cref{eq:mu_w} to get an intuitive sense of how the specification of $\Lambda$ in the CMSM affects the bounds on the causal estimands. 
When $\Lambda \to 1$, then the $(\Lambda^2 - 1)^{-1}$ term (and thus the denominator) in \Cref{eq:mu_w} tends to infinity.
As a result, the CAPO under $\Lambda$ converges to the empirical estimate of the CAPO — $\mu(w(\y); \x, \tf, \Lambda \to 1) \to \mut(\x, \tf)$ — as expected.
Thus, the supremum and infimum in \Cref{eq:CAPO_lower,eq:CAPO_upper} become independent of $w$, and the ignorance intervals concentrate on point estimates.
Next, consider complete relaxation of the ignorability assumption, $\Lambda \to \infty$. Then, the $(\Lambda^2 - 1)^{-1}$ term tends to zero, and we are left with,
\begin{equation*}
    \begin{split}
        \mu(w; \cdot, \Lambda \to \infty)
        \to \mut(\x, \tf) + \frac{\int_\Ycal w(\y) (\y - \mut(\x, \tf)) p(\y \mid \tf, \x)d\y}{\int_\Ycal w(\y)p(\y \mid \tf, \x)d\y,}
        = \mut(\x, \tf) + \!\!\!\! \E_{p(w(\y) \mid \x, \tf)} \!\!\!\! [\Y - \mut(\x, \tf)],
    \end{split}
\end{equation*}
where, $p(w(\y) \mid \x, \tf) \equiv \frac{w(\y) p(\y \mid \tf, \x)}{\int_\Ycal w(\y')p(\y' \mid \tf, \x)d\y'}$, a distribution over $\Y$ given $\X=x$ and $\Tf=\tf$.
Thus, when we \emph{relax} the ignorability assumption entirely, the CAPO can be anywhere in the range of $\Y$.

The parameter $\Lambda$ relates to the proportion of unexplained range in $\Y$ assumed to come from unobserved confounders after observing $\x$ and $\tf$. 
When a user sets $\Lambda$ to 1, they assume that the entire unexplained range of $\Y$ comes from unknown mechanisms independent of $\Tf$. 
As the user increases $\Lambda$, they attribute some of the unexplained range of $\Y$ to mechanisms causally connected to $\Tf$.
For bounded $\Yt$, this proportion can be calculated as:
\begin{equation*}
        \rho(\x, \tf; \Lambda) \coloneqq \frac{\overline{\mu}(\x, \tf; \Lambda) - \underline{\mu}(\x, \tf; \Lambda)}{\overline{\mu}(\x, \tf; \Lambda \to \infty) - \underline{\mu}(\x, \tf; \Lambda \to \infty)}
        = \frac{\overline{\mu}(\x, \tf; \Lambda) - \underline{\mu}(\x, \tf; \Lambda)}{\y_{\max} - \y_{\min} \mid \X=x, \Tf=\tf}.
\end{equation*}
The user can sweep over a set of $\Lambda$ values and report the bounds corresponding to a $\rho$ value they deem tolerable (e.g., $\rho=0.5$ yields bounds for the assumption that half the unexplained range in $Y$ is due to unobserved confounders). For unbounded outcomes, the limits can be estimated empirically by increasing $\Lambda$ to a large value. Refer to \Cref{fig:rho} in the appendix for a comparison between $\rho$ and $\Lambda$.

For another way to interpret $\Lambda$, in \Cref{app:kld} we $\Lambda$ can be presented as a bound on the Kullback–Leibler divergence between the nominal and complete propensity scores  through the relationship: $|\log{\left(\Lambda\right)}| \geq  D_{\mathrm{KL}}\left( P(\Yt \mid \Tf=\tf, \X=\x) || P(\Yt \mid \X=\x) \right)$.